\scanpage{5}
\textcircled{1.1.}
\[
\left\{\begin{aligned}
&k\ \text{corps}\\
&\bar k\ \text{clôture séparable},\quad \Gamma=\deleted{\unsure{\mathrm{Aut}_k|\mathrm{Aut}}}\mathrm{Gal}(\bar k/k)\quad\text{(groupe profini)}
\end{aligned}\right. \tag*{(1)}
\]
\[
\left\{\begin{aligned}
&\unsure{\mathscr X|\mathcal X|\mathfrak X}\ \text{schéma connexe sur }\bar k\\
&\unsure{\mathcal G|\mathscr G}=\mathrm{Aut}_{k}(\mathscr X)\\
&\unsure{\mathfrak G|G|\mathscr G}=\mathrm{Aut}_{\bar k}(\mathscr X)
\end{aligned}\right. \tag*{(2)}
\]
\margin{On suppose que $\bar k$ est séparablement clos dans $H^0(\mathscr X,\mathcal O_{\mathscr X})$ donc $\bar k$ est la clôture séparable de $k$ dans $H^{0}(\mathscr X,\mathcal O_{\unsure{\mathscr X|X|\mathfrak X}})$.}

\diagram{\fbox{\begin{tikzcd}[ampersand replacement=\&,column sep=small]
1\arrow[r]\&\mathfrak G\arrow[r]\&\mathcal G\arrow[r]\&\Gamma.
\end{tikzcd}}\quad(3)}

\textcircled{1.2} Si $k\subset k_1^{\unsure{\prime|\cdot}}\subset\bar k$, alors $k_1\longleftrightarrow\Gamma_1$ sous-groupe fermé de $\Gamma$

(4) $\boxed{\mathcal G_1^{\unsure{\prime|\cdot}}=\mathrm{Aut}_{k_1}(\mathscr X)}$ est alors l’image inverse de $\deleted{\unsure{\mathcal G}}\Gamma^{\unsure{\prime}}$ dans $\deleted{\unsure{\mathcal G}}\mathcal G$

\diagram{(5)\quad\begin{tikzcd}[ampersand replacement=\&,column sep=small]
1\arrow[r,dash]\&\mathfrak G\arrow[r,dash]\&\mathcal G\arrow[r]\&\Gamma\\
1\arrow[r,dash]\&\mathfrak G\arrow[u,equal]\arrow[r,dash]\&\mathcal G_1\arrow[u,hook]\arrow[r,"\text{cart}" description,yshift=0pt]\&\Gamma_1\arrow[u,hook]
\end{tikzcd}}

\textcircled{1.3} Une donnée de descente \deleted{\unsure{pour|de}} $\mathscr X$ de $\bar k$ à $k$ revient à un scindage de 3,

\diagram{(6)\quad\begin{tikzcd}[ampersand replacement=\&,column sep=small]
1\arrow[r]\&\mathfrak G\arrow[r,dash]\&\mathcal G\arrow[r]\&\Gamma\arrow[l,bend left=55]\arrow[r,"!"]\&1
\end{tikzcd}}

qui permet d’écrire
\[
\mathcal G\simeq\Gamma\mathbin{\underset{1/2}{\cdot}}\mathfrak G \tag*{(7)}
\]

\unsure{donc} sous réserve \inserted{de conditions} d’effectivité, \deleted{\unsure{il|et} \unsure{revient|revien|vraie}} (\unsure{OK|Ok} si $\mathscr X$ quasi-projectif sur $\bar k$) il revient au $=$ de se donner un couple $X/k$ et un iso
\[
X\otimes_k\bar k\simeq\mathscr X \tag*{(8)}
\]
% archive sheet number: 4
